\documentclass[11pt]{article}
\usepackage{amsmath,amssymb}
\usepackage{geometry}
\usepackage{enumitem}
\usepackage{multicol}
\geometry{margin=0.8in}
\setlength{\parindent}{0pt}
\setlength{\columnsep}{0.3in}

\begin{document}

\begin{center}
  {\Large\textbf{Functions and Average Rate of Change}}\\[3pt]
  {\large Algebra II Practice --- 45 Minutes}
\end{center}

Name: \rule{2.7in}{0.4pt}\hfill Date: \rule{1.5in}{0.4pt}

\section*{Instructions}
Complete all 12 problems in 45 minutes without a calculator. You may use this
cover sheet. Show key algebraic steps and give exact answers. Use interval
notation for domains and ranges. For conceptual questions, give a mathematical
reason or a counterexample; a yes/no answer alone is not sufficient.

\section*{Average Rate of Change (ARC)}
For $a<b$ and defined endpoint values $f(a)$ and $f(b)$,
\[
  \boxed{\operatorname{ARC}_{[a,b]}(f)=\frac{f(b)-f(a)}{b-a}}.
\]
This is the slope of the secant line through $(a,f(a))$ and $(b,f(b))$;
its units are output units per input unit. The difference quotient
$\dfrac{f(x+h)-f(x)}{h}$ requires $h\ne0$ and both function values to exist.

\section*{ARC for a Quadratic on $[a,b]$}
Let $f(x)=Ax^2+Bx+C$, where $A\ne0$. Since $b^2-a^2=(b-a)(a+b)$,
\[
  \operatorname{ARC}_{[a,b]}(f)
  =\frac{A(b^2-a^2)+B(b-a)}{b-a}
  =\frac{(b-a)\bigl(A(a+b)+B\bigr)}{b-a}.
\]
\[
  \boxed{\operatorname{ARC}_{[a,b]}(f)=A(a+b)+B}
  \qquad(a<b).
\]
In vertex form, $f(x)=A(x-h)^2+k$, the equivalent formula is
$\operatorname{ARC}_{[a,b]}(f)=A(a+b-2h)$.

\section*{Function Foundations}
\begin{multicols}{2}
\small
\textbf{Functions and notation}

A function assigns exactly one output to each input in its domain.
The domain is the set of allowed inputs; the range is the set of actual outputs.
The graph of a function passes the vertical line test. The notation $f(x)$
is a function value, not multiplication by $f$.

\medskip
\textbf{Domain restrictions}

Denominators cannot be zero. An even-root radicand must be nonnegative.
A real logarithm argument must be positive. Preserve restrictions from the
original expression even after algebraic simplification. For piecewise
functions, use the branch whose condition contains the input.

\columnbreak
\textbf{Composition}
\[
  (f\circ g)(x)=f(g(x)).
\]
Its domain requires $x\in\operatorname{Dom}(g)$ and
$g(x)\in\operatorname{Dom}(f)$.

\medskip
\textbf{One-to-one functions and inverses}

One-to-one means distinct inputs have distinct outputs; the graph passes the
horizontal line test. An inverse function requires a one-to-one original
function, possibly after restricting its domain.
\[
  f^{-1}(f(x))=x\quad\text{on }\operatorname{Dom}(f).
\]
Inverse functions exchange domain and range. The notation $f^{-1}$ means
inverse, not reciprocal.
\end{multicols}

\newpage
\begin{center}
  {\large\textbf{Functions and Average Rate of Change --- Practice}}
\end{center}

\begin{enumerate}[leftmargin=*,itemsep=0pt]
  \item Consider the relation
  \[
    R=\{(-3,4),\,(-1,2),\,(0,4),\,(2,-1)\}.
  \]
  State its domain and range. Determine whether $R$ is a function and
  whether its inverse relation is a function. Justify both decisions.
  \vspace{\stretch{1}}

  \item Let $f(x)=\dfrac{\sqrt{5-x}}{x-2}$. Find its natural real domain.
  Find $f(-4)$ and $f(5)$, and explain why $f(2)$ and $f(6)$ are not defined
  as real function values.
  \vspace{\stretch{1}}

  \item For $q(x)=-2x^2+7x-3$, find
  \[
    \operatorname{ARC}_{[-1,4]}(q),\qquad
    \operatorname{ARC}_{[1,3]}(q),\qquad
    \operatorname{ARC}_{[t,t+h]}(q)\quad(h>0).
  \]
  Give the last answer in terms of $t$ and $h$.
  \vspace*{\stretch{1}}
\end{enumerate}

\newpage
\begin{center}
  {\large\textbf{Functions and Average Rate of Change --- Continued}}
\end{center}

\begin{enumerate}[leftmargin=*,itemsep=0pt,resume]
  \item Let $f(x)=x^2-6x+8$. Find every interval $[a,b]$ of length $4$
  for which $\operatorname{ARC}_{[a,b]}(f)=0$. Is $f$ constant on the
  interval you found? Justify your answer using function values.
  \vspace{\stretch{1}}

  \item A quadratic $p(x)=Ax^2+Bx+C$ satisfies
  \[
    \operatorname{ARC}_{[0,2]}(p)=5,\qquad
    \operatorname{ARC}_{[1,4]}(p)=11,\qquad p(0)=-3.
  \]
  Determine $A$, $B$, and $C$.
  \vspace{\stretch{1}}

  \item The table gives the water remaining in a tank.
  \[
    \begin{array}{c|rrrr}
      t\text{ (minutes)} & 0 & 1 & 3 & 6\\ \hline
      V(t)\text{ (liters)} & 24 & 20 & 9 & 3
    \end{array}
  \]
  Find the average rates of change on $[0,3]$ and $[1,6]$, including units.
  A student says, ``The tank lost exactly 5 liters during each minute
  from $t=0$ to $t=3$.'' Does the table justify that claim? Explain.
  \vspace*{\stretch{1}}
\end{enumerate}

\newpage
\begin{center}
  {\large\textbf{Functions and Average Rate of Change --- Continued}}
\end{center}

\begin{enumerate}[leftmargin=*,itemsep=0pt,resume]
  \item Let
  \[
    p(x)=
    \begin{cases}
      x^2+1, & x<1,\\
      3x-1, & 1\le x<4,\\
      10-x, & x\ge4.
    \end{cases}
  \]
  Find $p(1)$, $p(4)$, $\operatorname{ARC}_{[-1,4]}(p)$, and
  $\operatorname{ARC}_{[1,4]}(p)$.
  \vspace{\stretch{1}}

  \item Let $f(x)=\sqrt{x-1}$ and $g(x)=x^2-3$.
  Find simplified formulas and domains for $f\circ g$ and $g\circ f$.
  Then find $\operatorname{ARC}_{[1,5]}(g\circ f)$.
  \vspace{\stretch{1}}

  \item Let $q(x)=(x-3)^2-4$. Explain why $q$ has no inverse function
  on its full real domain, using two specific inputs. Restrict its domain
  to $[3,\infty)$ and find $q^{-1}(x)$, together with the domain and range
  of that inverse.
  \vspace*{\stretch{1}}
\end{enumerate}

\newpage
\begin{center}
  {\large\textbf{Functions and Average Rate of Change --- Continued}}
\end{center}

\begin{enumerate}[leftmargin=*,itemsep=0pt,resume]
  \item Decide whether each statement is true or false. Give a reason
  for each true statement and a counterexample for each false statement.
  \begin{enumerate}[label=(\alph*),leftmargin=1.6em,itemsep=0.25em]
    \item If $\operatorname{ARC}_{[a,b]}(f)>0$, then $f$ must be strictly
    increasing throughout $[a,b]$.
    \item If two functions with domain $\mathbb R$ have equal ARC on every
    interval $[a,b]$, then they must be the same function.
    \item A function can have two distinct inputs with the same output.
    \item If both compositions exist, then $f\circ g=g\circ f$.
  \end{enumerate}
  \vspace{\stretch{1}}

  \item Consider
  \[
    f(x)=\frac{x^2-9}{x-3},\qquad g(x)=x+3,
  \]
  each with its natural real domain. Are these the same function?
  Explain. Compute their average rates of change on $[-1,2]$.
  Does agreement of those rates establish that the functions are equal?
  \vspace{\stretch{1}}

  \item For $f(x)=2x^2-3x+1$ and $h>0$, a student claims
  \[
    \frac{f(a+h)-f(a)}{h}=\frac{f(h)}{h}.
  \]
  Determine whether the claim holds for all real $a$ and positive $h$.
  Derive the correct simplified expression for the left side and give a
  specific counterexample if the claim is false. Also explain the
  difference between $f(a+h)$ and $f(a)+h$.
  \vspace*{\stretch{1}}
\end{enumerate}

\end{document}
